\chapter*{Introduzione} %Se si cambia il Titolo cambiare anche la riga successiva così che appia corretto nell'indice
\addcontentsline{toc}{chapter}{Introduzione} %Per far apparire Introduzione nell'indice (Il nome deve rispecchiare quello del chapter)
\pagenumbering{arabic} % Settaggio numerazione normale

Le vulnerabilità del software sono tra le principali cause di violazione della sicurezza informatica. Tra queste, la vulnerabilità \emph{format string} è nota da decenni eppure continua a comparire in codice reale: nasce da un errore tecnicamente semplice --- lasciare che l'utente controlli la stringa di formato passata a una funzione come \texttt{printf} --- ma le sue conseguenze possono arrivare fino a leggere informazioni sensibili e scrivere in modo arbitrario nella memoria del programma.

È durante CyberChallenge.IT --- il programma nazionale di addestramento in sicurezza informatica a cui ho preso parte --- che ho incontrato per la prima volta i fondamenti della sicurezza del software, un ambito centrale poiché una parte significativa degli attacchi informatici sfrutta proprio difetti di progettazione e vulnerabilità del codice. Delle vulnerabilità presentate, quella che ho trovato più preoccupante per la gravità delle sue possibili conseguenze è stata la format string.

Da questo interesse nasce il lavoro di tesi, che studia la vulnerabilità format string sull'architettura x86-64, attraverso la \texttt{printf} del linguaggio C, la funzione che meglio la rappresenta. Il contributo della tesi è duplice: da un lato spiega il meccanismo tecnico alla base e lo dimostra concretamente, sviluppando e sfruttando un programma vulnerabile scritto appositamente per lo studio, in un ambiente controllato e riproducibile; dall'altro valuta le difese dedicate alla format string, mettendone alla prova l'efficacia proprio sull'\emph{exploit} appena sviluppato.

Ne emerge che nessuna delle protezioni attive di default sul binario è sufficiente a impedirlo, e che perfino quelle create apposta per la format string ne bloccano solo una parte. L'unica difesa realmente efficace resta quindi alla radice: un uso corretto del linguaggio e delle sue funzioni, già in fase di scrittura del codice.

Il percorso che porta a questi risultati si sviluppa nei capitoli seguenti. Il Capitolo~\ref{chapter:background} introduce le nozioni tecniche necessarie a comprendere la vulnerabilità --- le funzioni variadiche del linguaggio C, il funzionamento della \texttt{printf} e la convenzione con cui i suoi argomenti vengono passati a basso livello sull'architettura x86-64 --- e fissa l'ambiente e gli strumenti usati nel resto del lavoro. Il Capitolo~\ref{chapter:fmt_vuln} presenta la vulnerabilità format string vera e propria, contrapponendo un uso corretto e uno vulnerabile della \texttt{printf} e mostrando, su due piccoli programmi di esempio, come da essa si possa sia leggere sia scrivere in memoria. Il Capitolo~\ref{chapter:exploit} combina queste tecniche in un caso pratico completo: dal punto di vista di un attaccante che dispone del solo binario compilato, senza codice sorgente, viene sviluppato un exploit che arriva fino all'ottenimento di una shell interattiva. Il Capitolo~\ref{chapter:mitigations} chiude la parte tecnica valutando le protezioni pensate specificamente per la format string, su due livelli complementari: uno a runtime, verificato sullo stesso exploit del Capitolo~\ref{chapter:exploit}, e uno in fase di compilazione, che permette di individuare la vulnerabilità alla radice.

Data la natura pratica delle tecniche discusse, va infine precisato che i contenuti presentati in questa tesi hanno esclusivamente finalità didattiche e di ricerca in ambito di sicurezza informatica: nessuna delle tecniche descritte è da intendersi come incentivo al loro utilizzo per scopi illeciti.
